%% 44장 직선 탐색과 준Newton 방법
\chapter{직선 탐색과 준Newton 방법}
\label{ch:line-search}

\abstract{다변수 함수를 최소화하는 방법은 대부분 내려갈 방향을 정하고 그 방향으로 얼마나 갈지 정하는 두 단계를 되풀이한다. 이 장은 둘째 단계부터 다룬다. 걸음이 함수값을 충분히 줄이면서도 너무 짧지 않아야 한다는 Wolfe 조건을 세우고 그런 걸음이 언제나 있다는 것과 방향과 기울기가 이루는 각이 직각에 한없이 가까워지지 않으면 그 조건만으로 기울기가 $0$으로 간다는 Zoutendijk의 정리를 증명한다. 방향으로는 Newton 방향과 그것을 함수값과 기울기만으로 흉내 내는 준Newton 방향을 다룬다. BFGS 갱신을 가중 Frobenius 노름에서 가장 작은 변화로 유도하고 양정치성을 지키는 데 필요한 조건을 Wolfe 조건의 곡률 조건이 보장한다는 것을 보인다. 초선형 수렴을 판정하는 Dennis--Moré 정리를 증명하고, Rosenbrock 함수에서 BFGS가 마지막 몇 걸음에 오차를 $10^{-6}$에서 $10^{-15}$로 줄이는 것을 확인한다. 이어서 행렬을 저장하지 않는 L-BFGS의 두 고리 재귀를 증명하고 상자 제약을 다루는 L-BFGS-B와 Hessian을 공액 기울기법으로 근사해 푸는 Newton-CG를 설명한다.}

\section{하강 방향과 직선 탐색}
\label{sec:line-search}

\subsection{하강 방향}

점 $x_k$에서 방향 $p_k$가 $\nabla f(x_k)^{\mathsf T}p_k<0$이면 $p_k$를 \term{하강 방향}{descent direction}이라 하는데 Taylor 전개 $f(x_k+\alpha p_k)=f(x_k)+\alpha\nabla f(x_k)^{\mathsf T}p_k+o(\alpha)$에서 보듯 이 방향으로 작은 $\alpha>0$만큼 가면 함수값이 줄기 때문이다. 음의 기울기 $-\nabla f(x_k)$와 방향 사이의 각을 $\theta_k$라 하면
\[
\cos\theta_k=\frac{-\nabla f(x_k)^{\mathsf T}p_k}{\|\nabla f(x_k)\|\,\|p_k\|}
\]
이고, 하강 방향은 $\cos\theta_k>0$인 방향이다. 42장의 경사 하강법은 $p_k=-\nabla f(x_k)$로 $\cos\theta_k=1$이다. 대칭 양정치 행렬 $B_k$로 $p_k=-B_k^{-1}\nabla f(x_k)$를 잡아도 $\nabla f^{\mathsf T}p_k=-\nabla f^{\mathsf T}B_k^{-1}\nabla f<0$이므로 하강 방향이 된다. Newton 방법은 $B_k=\nabla^2f(x_k)$이고 준Newton 방법은 $B_k$를 함수값과 기울기에서 만든다. 이하 $g_k=\nabla f(x_k)$로 쓰고, 방향을 정한 뒤의 한 변수 함수를
\[
\phi(\alpha)=f(x_k+\alpha p_k),\qquad\phi'(\alpha)=\nabla f(x_k+\alpha p_k)^{\mathsf T}p_k
\]
로 둔다. $\phi'(0)=g_k^{\mathsf T}p_k<0$이다.

\subsection{Wolfe 조건}

걸음 $\alpha_k$를 $\phi$의 최소점으로 정확히 잡을 필요는 없다. 43장의 방법으로 $\phi$를 정확히 최소화하려면 함수값을 수십 번 계산해야 하는데 방향부터가 어차피 근사이므로 그만한 정밀도는 낭비이다. 걸음은 함수값을 충분히 줄이면서 쓸데없이 짧지만 않으면 된다. 상수 $0<c_1<c_2<1$에 대해
\begin{align}
\phi(\alpha)&\le\phi(0)+c_1\alpha\phi'(0),\label{eq:armijo}\\
\phi'(\alpha)&\ge c_2\phi'(0)\label{eq:curvature}
\end{align}
를 \term{Wolfe 조건}{Wolfe conditions}이라 한다\cite{wolfe1969,wolfe1971}. \eqref{eq:armijo}는 \term{충분 감소 조건}{sufficient decrease condition} 또는 Armijo 조건으로\cite{armijo1966}, 실제 감소량이 접선이 예측하는 감소량의 $c_1$배 이상이라는 뜻이다. 충분 감소 조건만으로는 아주 작은 $\alpha$도 통과하므로 \eqref{eq:curvature}의 \term{곡률 조건}{curvature condition}을 더한다. $\phi'(\alpha)$가 여전히 크게 음수이면 더 가서 더 줄일 수 있으므로, 기울기가 처음의 $c_2$배만큼은 평평해진 곳까지 가라는 조건이다. 곡률 조건을 $|\phi'(\alpha)|\le c_2|\phi'(0)|$로 바꾼 것을 \term{강한 Wolfe 조건}{strong Wolfe conditions}이라 한다. 강한 조건은 $\phi'$이 크게 양수인 곳, 곧 최소점을 한참 지나친 곳까지 막는다. 준Newton 방법에서는 보통 $c_1=10^{-4}$, $c_2=0.9$를 쓴다.

\begin{proposition}
\label{prop:wolfe-exists}
$f$가 연속 미분 가능하고, $p_k$가 하강 방향이며, $\phi$가 $\alpha>0$에서 아래로 유계라 하자. $0<c_1<c_2<1$이면 강한 Wolfe 조건을 만족하는 $\alpha$들로 이루어진 구간이 있다.
\end{proposition}

\begin{proof}
직선 $\ell(\alpha)=\phi(0)+c_1\alpha\phi'(0)$은 기울기가 음수이므로 아래로 유계가 아니고, $\phi$는 아래로 유계이므로 둘은 $\alpha>0$에서 만난다. 또 $\phi'(0)<c_1\phi'(0)$이므로 작은 $\alpha>0$에서 $\phi<\ell$이다. 만나는 점 가운데 가장 작은 것을 $\alpha'$이라 하면 $(0,\alpha')$에서 $\phi<\ell$이고 $\phi(\alpha')=\ell(\alpha')$이다. 평균값 정리로 어떤 $\alpha''\in(0,\alpha')$에서
\[
\phi'(\alpha'')=\frac{\phi(\alpha')-\phi(0)}{\alpha'}=c_1\phi'(0)
\]
이다. $\alpha''<\alpha'$이므로 \eqref{eq:armijo}를 엄격한 부등호로 만족하고, $|\phi'(\alpha'')|=c_1|\phi'(0)|<c_2|\phi'(0)|$이므로 강한 곡률 조건도 엄격한 부등호로 만족한다. 두 부등식의 양변이 연속이므로 $\alpha''$를 품는 어떤 열린 구간 전체가 조건을 만족한다. 다만 $\phi$가 여러 번 오르내리면 조건을 만족하는 걸음이 구간 여러 개로 나뉠 수 있다.
\end{proof}

\subsection{Zoutendijk의 정리}

Wolfe 조건이 왜 알맞은지는 다음 정리가 보여 준다. 방향이 기울기와 직각에 한없이 가까워지지만 않으면 Wolfe 걸음만으로 기울기가 $0$으로 간다.

\begin{theorem}[Zoutendijk]
\label{thm:zoutendijk}
$f$가 아래로 유계이고 초기 준위 집합 $\{x:f(x)\le f(x_0)\}$의 열린 근방에서 연속 미분 가능하며, 그곳에서 $\nabla f$가 상수 $L$로 Lipschitz 연속이라 하자. 모든 $p_k$가 하강 방향이고 모든 $\alpha_k$가 Wolfe 조건을 만족하면
\[
\sum_{k\ge0}\cos^2\theta_k\,\|g_k\|^2<\infty
\]
이다. 특히 어떤 $\delta>0$에 대해 언제나 $\cos\theta_k\ge\delta$이면 $\|g_k\|\to0$이다.
\end{theorem}

\begin{proof}
곡률 조건 \eqref{eq:curvature}에서 $g_{k+1}^{\mathsf T}p_k\ge c_2g_k^{\mathsf T}p_k$이므로 $(g_{k+1}-g_k)^{\mathsf T}p_k\ge(c_2-1)g_k^{\mathsf T}p_k$이다. 한편 Lipschitz 조건과 Cauchy--Schwarz 부등식으로 $(g_{k+1}-g_k)^{\mathsf T}p_k\le\alpha_kL\|p_k\|^2$이다. 두 식에서
\[
\alpha_k\ge\frac{(c_2-1)g_k^{\mathsf T}p_k}{L\|p_k\|^2}
\]
이다. 이 하한을 충분 감소 조건에 넣으면 $g_k^{\mathsf T}p_k<0$이므로
\begin{align*}
f(x_{k+1})&\le f(x_k)+c_1\alpha_kg_k^{\mathsf T}p_k\le f(x_k)-\frac{c_1(1-c_2)}L\,\frac{(g_k^{\mathsf T}p_k)^2}{\|p_k\|^2}\\
&=f(x_k)-\frac{c_1(1-c_2)}L\cos^2\theta_k\|g_k\|^2
\end{align*}
이다. 함수값이 줄어들므로 모든 $x_k$가 준위 집합에 있고, 이 부등식을 $k=0,\dots,K$에 대해 더하면 $\sum_{k\le K}\cos^2\theta_k\|g_k\|^2\le L(f(x_0)-\inf f)/(c_1(1-c_2))$이다. $K\to\infty$로 보내면 된다. 둘째 주장은 수렴하는 급수의 항이 $0$으로 가는 데서 나온다.
\end{proof}

이 정리는 수렴의 속도는 말하지 않는다. 경사 하강법은 $\cos\theta_k=1$이어서 $g_k\to0$이 보장되지만, 42장에서 보았듯이 조건수가 크면 아주 느리다. 2차원 Rosenbrock 함수 $f(x,y)=100(y-x^2)^2+(1-x)^2$에서 $(-1.2,1)$로 시작해 강한 Wolfe 걸음($c_1=10^{-4}$, $c_2=0.9$)의 경사 하강법을 $\|g\|_\infty\le10^{-5}$까지 돌리면 $7{,}095$번 반복하고 함수값을 $56{,}219$번 계산한다. 같은 판정에서 BFGS는 $32$번 반복하고 함수값을 $39$번 계산하니, 걸음보다 방향을 잘 고르는 쪽이 훨씬 중요하다. 그래도 Zoutendijk의 정리\cite{nocedal2006}가 있으면 방향을 바꾸면서도 수렴을 지킬 수 있다. 어떤 방법이든 가끔 경사 하강 걸음을 섞으면 그 걸음들에서 $\cos\theta_k=1$이므로 기울기가 $0$으로 가는 부분열이 생긴다.

\subsection{걸음을 찾는 알고리즘}

Wolfe 조건을 만족하는 걸음을 찾는 표준 알고리즘은 Moré와 Thuente의 것으로\cite{more1994linesearch}, MINPACK-2의 \code{dcsrch}로 널리 퍼졌다. 이 알고리즘은 43장에서 최솟값을 낀 구간을 유지했듯이 조건을 만족하는 걸음을 낀 구간을 유지한다. 시험 걸음마다 함수값과 도함수를 계산해 구간을 줄이고, 다음 시험 걸음은 구간 끝의 함수값과 도함수로 만든 3차 또는 2차 보간의 최소점에서 고른다. 보간이 구간을 충분히 줄이지 못하면 이분하는 안전 장치가 있다. 구간이 아직 없을 때는 걸음을 키워 가며 구간을 찾는다. 보간에 도함수를 쓰므로 함수값만 쓰는 43장의 방법보다 훨씬 빨리 끝나고, 준Newton 방법에서는 첫 시험 걸음 $\alpha=1$이 대부분 그대로 받아들여진다.

scipy와 stardust의 BFGS와 Newton-CG는 먼저 이 \code{dcsrch}로 강한 Wolfe 걸음을 찾고, 실패하면 Nocedal과 Wright의 책에 있는 다른 알고리즘(걸음을 두 배씩 키우는 단계와 3차·2차 보간으로 구간을 줄이는 zoom 단계)으로 다시 찾는다\cite{nocedal2006}. 둘 다 실패하면 방법이 멈추고 그 사실을 알린다. 첫 시험 걸음은 직전 반복의 함수값 감소에서 어림한다. 이번 반복도 직전 반복만큼 줄인다고 보면 2차 보간의 최소점이 $\alpha=2(f_k-f_{k-1})/\phi'(0)$이므로, 이것에 $1.01$을 곱한 값과 $1$ 가운데 작은 것으로 시작한다. 첫 반복에서 BFGS는 직전 값 대신 $f_{-1}=f_0+\|g_0\|/2$로 두는데, 그러면 $p_0=-g_0$에서 첫 시험 걸음이 $\alpha=\min(1,1.01/\|g_0\|)$이 되어, 첫 걸음의 길이 $\alpha\|g_0\|$는 $\min(\|g_0\|,1.01)$이다. 기울기가 크면 길이가 대략 $1$로 묶인다. Newton-CG는 첫 반복에서 $\alpha=1$로 시작한다.

\section{Newton 방법}
\label{sec:newton-min}

$f$를 $x_k$ 둘레의 이차 모형 $m_k(p)=f(x_k)+g_k^{\mathsf T}p+p^{\mathsf T}\nabla^2f(x_k)p/2$로 근사하고 그 최소점으로 가는 것이 Newton 방법이다. $\nabla^2f(x_k)$가 양정치이면 모형의 최소점은
\[
p_k=-\nabla^2f(x_k)^{-1}g_k
\]
이다. 이것은 $\nabla f(x)=0$에 34장의 Newton 방법을 적용한 것과 같으므로, $\nabla^2f$가 최소점 $x_\ast$ 근방에서 Lipschitz 연속이고 $\nabla^2f(x_\ast)$가 양정치이면 $x_\ast$ 가까이에서 이차로 수렴한다. 또 Newton 방향은 좌표를 선형으로 바꾸어도 같은 걸음을 준다. $x=Ay$로 두면 $y$에 대한 기울기는 $A^{\mathsf T}g$, Hessian은 $A^{\mathsf T}\nabla^2fA$이므로 $y$의 Newton 걸음은 $A^{-1}p_k$이다. 경사 하강법에는 없는 이 성질 때문에 Newton 방법은 조건수에 둔감하다.

최소점에서 먼 곳에서는 사정이 다르다. $\nabla^2f(x_k)$가 양정치가 아니면 $p_k$는 하강 방향이 아닐 수 있고, 모형이 아래로 유계가 아니면 최소점이 없다. 직선 탐색과 함께 쓰는 Newton 방법은 이때 $\nabla^2f(x_k)+\tau I$처럼 행렬을 고쳐 양정치로 만든다. Rosenbrock 함수를 $(-1.2,1)$에서 시작하면 이 경로에서는 Hessian이 줄곧 양정치이고, Wolfe 걸음의 Newton 방법은 $21$번 반복으로 $\|g\|_\infty\le10^{-5}$에 이른다. 고치는 방법은 다음 장의 신뢰 영역이 더 깔끔하게 다룬다. 더 근본적인 걸림돌은 비용으로, Hessian을 계산하고 저장하는 데 $n^2$, 분해하는 데 $n^3$이 든다. 다음 세 절은 이 비용을 피하는 두 길, 곧 Hessian을 기울기로부터 쌓아 가는 준Newton 방법과 Hessian과 벡터의 곱만 쓰는 Newton-CG이다.

\section{BFGS 방법}
\label{sec:bfgs}

\subsection{할선 조건}

준Newton 방법은 $\nabla^2f(x_k)$ 대신 근사 $B_k$를 쓰되, 걸음마다 얻는 기울기의 변화로 그것을 고친다. 걸음 $s_k=x_{k+1}-x_k$와 기울기의 변화 $y_k=g_{k+1}-g_k$ 사이에는 평균값 정리로
\[
y_k=\Bigl(\int_0^1\nabla^2f(x_k+ts_k)\,dt\Bigr)s_k
\]
가 성립하므로, 새 근사도 이 관계를 만족하기를 요구한다. 곧
\[
B_{k+1}s_k=y_k
\]
이고, 이것을 \term{할선 조건}{secant condition}이라 한다. 34장의 Broyden 방법이 방정식계에서 쓴 조건과 같다. 다만 여기서는 $B_{k+1}$이 대칭이고 양정치이기를 바란다. 양정치인 $B_{k+1}$이 할선 조건을 만족하면 $s_k^{\mathsf T}y_k=s_k^{\mathsf T}B_{k+1}s_k>0$이어야 하므로
\[
s_k^{\mathsf T}y_k>0
\]
이 필요하다. 이것을 \term{곡률 조건}{curvature condition}이라 하는데 Wolfe 조건의 곡률 조건이 바로 이 부등식을 보장한다. \eqref{eq:curvature}에서 $g_{k+1}^{\mathsf T}s_k\ge c_2g_k^{\mathsf T}s_k$이므로
\[
s_k^{\mathsf T}y_k\ge(c_2-1)\,g_k^{\mathsf T}s_k=(1-c_2)\,\alpha_k\,|g_k^{\mathsf T}p_k|>0
\]
이다. 준Newton 방법이 Wolfe 조건과 짝을 이루는 것은 볼록이 아닌 함수에서도 Wolfe 걸음이면 곡률 조건이 성립하기 때문이다.

방정식을 실제로 풀지 않으려면 $B_k$ 대신 그 역 $H_k=B_k^{-1}$을 유지하는 편이 낫다. 그러면 방향이 행렬과 벡터의 곱 $p_k=-H_kg_k$이고, 할선 조건은 $H_{k+1}y_k=s_k$이다.

\subsection{가장 작은 변화}

$n>1$이면 할선 조건은 $H_{k+1}$의 $n(n+1)/2$개 자유도 가운데 $n$개만 정하므로, 나머지는 $H_k$에서 가장 적게 바꾸는 것으로 정한다. 변화의 크기를 재는 노름에 따라 다른 갱신이 나오는데, 다음 정리의 가중 노름이 BFGS 갱신을 준다. 이 정리에서 무게 행렬 $W$는 $Ws_k=y_k$를 만족하는 아무 대칭 양정치 행렬이다. 예컨대 위의 평균 Hessian이 양정치이면 그것이 그런 행렬이다. 가중 노름 $\|A\|_W=\|W^{1/2}AW^{1/2}\|_F$는 문제의 척도를 바꾸어도 바뀌지 않는다.

\begin{theorem}
\label{thm:bfgs-variational}
$H_k$가 대칭이고 $s_k^{\mathsf T}y_k>0$이라 하자. $Ws_k=y_k$인 대칭 양정치 행렬 $W$에 대해 문제
\[
\min_H\|H-H_k\|_W\quad\text{조건: }H=H^{\mathsf T},\ Hy_k=s_k
\]
의 유일한 해는
\begin{equation}
H_{k+1}=(I-\rho_ks_ky_k^{\mathsf T})\,H_k\,(I-\rho_ky_ks_k^{\mathsf T})+\rho_ks_ks_k^{\mathsf T},\qquad\rho_k=\frac1{y_k^{\mathsf T}s_k}
\label{eq:bfgs}
\end{equation}
이다.
\end{theorem}

\begin{proof}
첨자 $k$를 떼고, $\hat H=W^{1/2}HW^{1/2}$, $\hat H_k=W^{1/2}H_kW^{1/2}$, $\hat y=W^{-1/2}y$로 바꾸어 쓰자. $Ws=y$에서 $W^{1/2}s=W^{-1/2}y=\hat y$이므로 조건 $Hy=s$는 $\hat H\hat y=\hat y$가 되고, 목적은 $\|\hat H-\hat H_k\|_F$이다. $E=\hat H-\hat H_k$로 두면 대칭 행렬 $E$가 $E\hat y=r$, $r=\hat y-\hat H_k\hat y$를 만족하는 가운데 Frobenius 노름이 가장 작은 것을 찾는 문제이다.

$q=\hat y/\|\hat y\|$, $t=r/\|\hat y\|$로 두고 $q$를 첫 벡터로 하는 정규직교 기저를 잡으면, 조건은 $E$의 첫 열이 $t$라는 것이다. 대칭이므로 첫 행도 $t^{\mathsf T}$로 정해지고, 나머지 $(n-1)\times(n-1)$ 블록 $C$만 자유롭다. $\|E\|_F^2=(q^{\mathsf T}t)^2+2\|t-(q^{\mathsf T}t)q\|^2+\|C\|_F^2$이므로 $C=0$일 때만 최소이고, 그 해는
\[
E=tq^{\mathsf T}+qt^{\mathsf T}-(q^{\mathsf T}t)\,qq^{\mathsf T}=\frac{r\hat y^{\mathsf T}+\hat yr^{\mathsf T}}{\hat y^{\mathsf T}\hat y}-\frac{(\hat y^{\mathsf T}r)\,\hat y\hat y^{\mathsf T}}{(\hat y^{\mathsf T}\hat y)^2}
\]
이다. $P=\hat y\hat y^{\mathsf T}/(\hat y^{\mathsf T}\hat y)$로 쓰면 $r\hat y^{\mathsf T}/(\hat y^{\mathsf T}\hat y)=P-\hat H_kP$, $\hat yr^{\mathsf T}/(\hat y^{\mathsf T}\hat y)=P-P\hat H_k$, $(\hat y^{\mathsf T}r)\hat y\hat y^{\mathsf T}/(\hat y^{\mathsf T}\hat y)^2=P-P\hat H_kP$이므로
\[
\hat H=\hat H_k+E=\hat H_k-\hat H_kP-P\hat H_k+P\hat H_kP+P=(I-P)\hat H_k(I-P)+P
\]
이다. $\hat y^{\mathsf T}\hat y=y^{\mathsf T}W^{-1}y=y^{\mathsf T}s=1/\rho$이고 $W^{-1/2}PW^{1/2}=\rho\,sy^{\mathsf T}$, $W^{-1/2}PW^{-1/2}=\rho\,ss^{\mathsf T}$이므로, 양쪽에 $W^{-1/2}$를 곱해 되돌리면 \eqref{eq:bfgs}가 나온다.
\end{proof}

\eqref{eq:bfgs}의 $H_{k+1}$은 $W$에 의존하지 않으므로, 할선 조건을 만족하는 어느 무게로 재든 같은 갱신이 가장 작은 변화이다. 이 갱신은 1970년에 Broyden, Fletcher, Goldfarb, Shanno가 서로 다른 길로 같은 식에 이르러 네 사람의 머리글자로 \term{BFGS 갱신}{BFGS update}이라 불린다\cite{broyden1970,fletcher1970,goldfarb1970,shanno1970}. 같은 문제를 $B$에 대해 세워 $\|W^{-1/2}(B-B_k)W^{-1/2}\|_F$를 최소화하면 \term{DFP 갱신}{DFP update}이 나오는데, 이것은 BFGS 갱신의 $H$와 $B$, $s$와 $y$를 맞바꾼 것이다. 실제로는 BFGS가 훨씬 나은데, 근사가 나빠졌을 때 스스로 고치는 성질이 강하기 때문으로 알려져 있다\cite{nocedal2006}.

\begin{proposition}
\label{prop:bfgs-pd}
$H_k$가 대칭 양정치이고 $s_k^{\mathsf T}y_k>0$이면 \eqref{eq:bfgs}의 $H_{k+1}$도 대칭 양정치이다.
\end{proposition}

\begin{proof}
대칭은 식에서 바로 보인다. $z\ne0$에 대해 $w=z-\rho_k(s_k^{\mathsf T}z)\,y_k$로 두면 $(I-\rho_ky_ks_k^{\mathsf T})z=w$이므로
\[
z^{\mathsf T}H_{k+1}z=w^{\mathsf T}H_kw+\rho_k(s_k^{\mathsf T}z)^2
\]
이고, 두 항 모두 음이 아니다. 둘째 항이 $0$이면 $s_k^{\mathsf T}z=0$이므로 $w=z\ne0$이고, 그러면 첫째 항이 양수이다.
\end{proof}

\cref{alg:bfgs}가 BFGS 방법이다. 처음 $H_0$는 보통 단위 행렬이다. 행렬과 벡터의 곱과 갱신이 반복마다 $O(n^2)$이므로, Hessian을 분해하는 Newton 방법의 $O(n^3)$보다 싸다.

\begin{algorithm}[t]
\begin{algorithmic}[1]
\Require 함수 $f$와 기울기 $\nabla f$, 시작점 $x_0$, 허용오차 $\mathrm{gtol}$
\Ensure 근사 최소점 $x_k$
\State $H_0\gets I$, $g_0\gets\nabla f(x_0)$, $k\gets0$
\While{$\|g_k\|_\infty>\mathrm{gtol}$}
  \State $p_k\gets-H_kg_k$
  \State 강한 Wolfe 조건을 만족하는 $\alpha_k$를 찾는다 \Comment{$c_1=10^{-4}$, $c_2=0.9$}
  \State $s_k\gets\alpha_kp_k$, $x_{k+1}\gets x_k+s_k$, $g_{k+1}\gets\nabla f(x_{k+1})$, $y_k\gets g_{k+1}-g_k$
  \State $\rho_k\gets1/(y_k^{\mathsf T}s_k)$, $v\gets H_ky_k$
  \State $H_{k+1}\gets H_k-\rho_k(s_kv^{\mathsf T}+vs_k^{\mathsf T})+(\rho_k^2\,y_k^{\mathsf T}v+\rho_k)\,s_ks_k^{\mathsf T}$ \Comment{\eqref{eq:bfgs}를 펼친 꼴}
  \State $k\gets k+1$
\EndWhile
\end{algorithmic}
\caption{BFGS 방법}
\label{alg:bfgs}
\end{algorithm}

\subsection{초선형 수렴}

준Newton 방법의 $B_k$는 $\nabla^2f(x_\ast)$로 수렴하지 않을 수도 있다. 그런데도 Newton 방법에 가까운 속도를 내는 까닭을 다음 정리가 설명하는데 $B_k$는 걸음의 방향에서만 Hessian과 같으면 된다.

\begin{theorem}[Dennis--Moré]
\label{thm:dennis-more}
$f$가 두 번 연속 미분 가능하고 $\nabla^2f$가 $x_\ast$의 근방에서 Lipschitz 연속이며, $\nabla f(x_\ast)=0$이고 $\nabla^2f(x_\ast)$가 양정치라 하자. 반복 $x_{k+1}=x_k+p_k$, $B_kp_k=-g_k$가 $x_\ast$로 수렴하고 모든 $k$에서 $x_k\ne x_\ast$라 하자. 그러면 수렴이 초선형, 곧 $\|x_{k+1}-x_\ast\|/\|x_k-x_\ast\|\to0$인 것은
\[
\lim_{k\to\infty}\frac{\|(B_k-\nabla^2f(x_\ast))\,p_k\|}{\|p_k\|}=0
\]
인 것과 동치이다.
\end{theorem}

\begin{proof}
$e_k=x_k-x_\ast$, $G=\nabla^2f(x_\ast)$로 쓰자. 적분 꼴의 평균값 정리와 Lipschitz 조건으로
\[
g_{k+1}=g_k+\int_0^1\nabla^2f(x_k+tp_k)\,p_k\,dt=g_k+Gp_k+O\bigl((\|e_k\|+\|e_{k+1}\|)\,\|p_k\|\bigr)
\]
이다. $g_k=-B_kp_k$를 넣고 $e_k\to0$을 쓰면
\[
g_{k+1}=-(B_k-G)\,p_k+o(\|p_k\|)
\]
이다. 그러므로 정리의 조건은 $\|g_{k+1}\|/\|p_k\|\to0$과 동치이다. 한편 $G$가 양정치이므로 $x_\ast$ 근방에서 $\|g_{k+1}\|$은 $\|e_{k+1}\|$과 비교 가능하다. 곧 어떤 $0<m\le M$에 대해 $m\|e_{k+1}\|\le\|g_{k+1}\|\le M\|e_{k+1}\|$이다.

조건이 성립한다고 하자. $\|p_k\|\le\|e_k\|+\|e_{k+1}\|$이므로 $\|e_{k+1}\|\le\|g_{k+1}\|/m=o(\|e_k\|+\|e_{k+1}\|)$이고, 이것을 $\|e_{k+1}\|$에 대해 풀면 $\|e_{k+1}\|=o(\|e_k\|)$이다. 거꾸로 초선형이면 $\|p_k\|\ge\|e_k\|-\|e_{k+1}\|=\|e_k\|(1-o(1))$이므로 $\|g_{k+1}\|\le M\|e_{k+1}\|=o(\|e_k\|)=o(\|p_k\|)$이다.
\end{proof}

이 정리는 Dennis와 Moré가 1974년에 증명했다\cite{dennis1974}. 직선 탐색과 함께 쓰려면 여기에 더해 끝에서는 걸음 $\alpha_k=1$이 받아들여져야 한다. 같은 논문은 정리의 조건이 성립하고 $c_1<1/2$이면 충분히 큰 $k$에서 $\alpha_k=1$이 Wolfe 조건을 만족함을 보였다. 그래서 직선 탐색은 언제나 $\alpha=1$을 먼저 시험해야 한다. BFGS가 준위 집합에서 강볼록이고(곧 Hessian이 균등하게 양정치이고) Hessian이 Lipschitz 연속인 함수에서 이 조건을 만족하여 초선형으로 수렴한다는 것은 Powell(1976)과 Byrd--Nocedal(1989)의 결과이다\cite{nocedal2006}.

Rosenbrock 함수에서 확인해 보자. \cref{alg:bfgs}를 scipy의 \code{line\_search}와 함께 직접 짜서 $\|g\|_\infty\le10^{-10}$까지 돌리면, $(-1.2,1)$로 시작한 BFGS는 처음 스무 걸음쯤 굽은 골짜기를 따라 기어가며 걸음 길이도 $1$보다 작은 것이 섞인다. $27$번째 반복부터는 모든 걸음이 $\alpha=1$이고, 최소점까지의 거리가 $3.3\times10^{-3}$, $6.2\times10^{-5}$, $1.4\times10^{-6}$, $1.4\times10^{-8}$, $1.4\times10^{-10}$, $6.3\times10^{-15}$으로 줄어, 이웃한 비가 $2\times10^{-2}$ 안팎에서 마지막에 $5\times10^{-5}$까지 떨어진다. 같은 걸음들에서 Dennis--Moré의 비 $\|(B_k-G)p_k\|/\|p_k\|$는 $\|G\|_2\approx1002$에 견주어 $3.6$, $1.7$, $4.5$, $7.3\times10^{-2}$, $2.0\times10^{-2}$, $2.5\times10^{-3}$으로 흔들리며 준다.

\subsection{유한 정밀도에서의 갱신}

\eqref{eq:bfgs}를 그대로 계산하면 $n\times n$ 행렬곱이 두 번이어서 $O(n^3)$이 든다. scipy는 이 꼴로 계산한다. stardust는 같은 식을 펼친
\[
H_{k+1}=H_k-\rho_k(s_kv^{\mathsf T}+vs_k^{\mathsf T})+(\rho_k^2\,y_k^{\mathsf T}v+\rho_k)\,s_ks_k^{\mathsf T},\qquad v=H_ky_k
\]
로 $O(n^2)$에 계산한다. 수학적으로 같은 식이지만 반올림은 달라서, 극단적인 경우인 $n=1$에서 펼친 식은 $H-2H+H+s/y$를 계산하는데 $s/y$가 $H$보다 아주 작으면 앞의 세 항이 상쇄되며 생긴 반올림 오차가 $s/y$를 삼킨다. $H=1$, $s=10^{-9}$, $y=10^8$이면 참값 $s/y=10^{-17}$ 대신 계산 순서에 따라 $0$이나 $2.2\times10^{-16}$이 나온다. 곱의 꼴 $(1-\rho sy)H(1-\rho ys)+\rho s^2$은 앞의 괄호가 정확히 $0$이 되어 $10^{-17}$을 준다. 그래서 stardust는 $n=1$에서만 곱의 꼴을 쓴다. 9장의 상쇄가 여기서도 같은 모습으로 나타난다. 펼친 식의 $(i,j)$와 $(j,i)$ 원소는 같은 항을 순서만 바꾼 식이므로, 같은 순서로 계산하면 $H_{k+1}$은 비트까지 대칭으로 남는다.

\begin{aside}[Davidon의 보고서]
최초의 준Newton 방법은 아르곤 국립연구소의 물리학자 William Davidon이 1959년에 쓴 보고서에 있었다. 긴 계산이 끝나기 전에 컴퓨터가 자주 멈추는 것에 지쳐 더 빠른 방법을 찾은 것이라고 한다. 그 보고서는 학술지에 실리지 못했지만 Fletcher와 Powell이 1963년에 그 방법을 다듬어 알렸고\cite{fletcher1963}, 이것이 DFP 방법이다. 보고서는 32년 뒤인 1991년에 SIAM Journal on Optimization의 첫 권 첫 호 첫 논문으로 실렸다\cite{davidon1991}. Davidon은 반전 운동가이기도 해서, 1971년 펜실베이니아주 미디어의 FBI 사무소에 침입해 불법 사찰 문서를 폭로한 사람들 가운데 하나였다는 사실이 2014년에야 알려졌다.
\end{aside}

\section{제한 메모리 BFGS}
\label{sec:lbfgs}

\subsection{두 고리 재귀}

$n$이 수만이면 $n\times n$ 행렬 $H_k$를 저장할 수도 없다. 그런데 \eqref{eq:bfgs}는 $H_k$를 쌍 $(s_k,y_k)$로 고친 것이므로, 최근 $m$개의 쌍과 처음 행렬 $H_k^0$만 있으면 $H_k$를 저장하지 않고도 $H_kg_k$를 계산할 수 있다. $V_i=I-\rho_iy_is_i^{\mathsf T}$로 쓰면 \eqref{eq:bfgs}는 $H_{i+1}=V_i^{\mathsf T}H_iV_i+\rho_is_is_i^{\mathsf T}$이다. 이 재귀를 $i=k-m$에서 $H_{k-m}=H_k^0$으로 시작해 $k-1$까지 쓴 것이 제한 메모리의 근사 $H_k$이다. \cref{alg:two-loop}는 이 $H_kq$를 $O(mn)$에 계산한다\cite{nocedal1980}.

\begin{algorithm}[t]
\begin{algorithmic}[1]
\Require 벡터 $q$, 쌍 $(s_i,y_i)$와 $\rho_i=1/(y_i^{\mathsf T}s_i)$ ($i=k-m,\dots,k-1$), 처음 행렬 $H_k^0$
\Ensure $r=H_kq$
\For{$i=k-1,k-2,\dots,k-m$} \Comment{첫 고리: 새것에서 옛것으로}
  \State $\alpha_i\gets\rho_is_i^{\mathsf T}q$, $q\gets q-\alpha_iy_i$
\EndFor
\State $r\gets H_k^0q$
\For{$i=k-m,\dots,k-1$} \Comment{둘째 고리: 옛것에서 새것으로}
  \State $\beta\gets\rho_iy_i^{\mathsf T}r$, $r\gets r+(\alpha_i-\beta)\,s_i$
\EndFor
\State \Return $r$
\end{algorithmic}
\caption{L-BFGS의 두 고리 재귀}
\label{alg:two-loop}
\end{algorithm}

\begin{proposition}
\cref{alg:two-loop}는 $H_kq$를 돌려준다.
\end{proposition}

\begin{proof}
첫 고리가 $i$에서 만드는 벡터를 $q_i$라 하자. 처음 $q_k=q$이고 $q_i=q_{i+1}-\alpha_iy_i$, $\alpha_i=\rho_is_i^{\mathsf T}q_{i+1}$이므로 $q_i=V_iq_{i+1}$이다. 둘째 고리에서 $i$번째 갱신 직전의 $r$을 $r_i$라 하면 $r_{k-m}=H_k^0q_{k-m}=H_{k-m}q_{k-m}$이고, $\beta=\rho_iy_i^{\mathsf T}r_i$이므로
\[
r_{i+1}=r_i-\rho_is_iy_i^{\mathsf T}r_i+\alpha_is_i=V_i^{\mathsf T}r_i+\alpha_is_i
\]
이다. $r_i=H_iq_i$라 가정하면
\[
H_{i+1}q_{i+1}=V_i^{\mathsf T}H_iV_iq_{i+1}+\rho_is_is_i^{\mathsf T}q_{i+1}=V_i^{\mathsf T}H_iq_i+\alpha_is_i=r_{i+1}
\]
이므로 귀납법으로 $r_k=H_kq_k=H_kq$이다.
\end{proof}

처음 행렬은 보통 $H_k^0=\gamma_kI$, $\gamma_k=s_{k-1}^{\mathsf T}y_{k-1}/(y_{k-1}^{\mathsf T}y_{k-1})$로 잡는다. $\gamma_k$는 최근 걸음 방향에서 본 Hessian의 역의 크기를 어림하므로 첫 시험 걸음 $\alpha=1$이 대개 알맞은 길이가 된다. 메모리 $m$은 $3$에서 $20$ 사이면 충분하다\cite{liu1989}. 100차원 Rosenbrock 함수(scipy의 \code{rosen}, 이웃한 좌표를 잇는 꼴)를 $(-1.2,1,-1.2,1,\dots)$에서 L-BFGS-B로 풀면 반복 수가 $m=1,3,5,10,20,50$에서 $763$, $542$, $532$, $514$, $507$, $489$이다. $m$을 $3$에서 $50$으로 늘려도 반복은 $10\%$쯤만 준다. 같은 문제에서 행렬 전체를 쓰는 BFGS는 $538$번 반복한다.

\subsection{상자 제약과 L-BFGS-B}

분산은 음수일 수 없고 확률은 $[0,1]$에 있어야 하듯이, 많은 문제에서 변수에는 상한과 하한이 있다. Byrd, Lu, Nocedal, Zhu의 \term{L-BFGS-B}{L-BFGS-B}는 제약 $l\le x\le u$ 아래에서 L-BFGS를 쓰는 방법이다\cite{byrd1995lbfgsb,zhu1997lbfgsb}. 반복마다 세 단계를 거친다.

첫째는 \term{일반화 Cauchy 점}{generalized Cauchy point}이다. 상자로의 사영 $P(x)$(성분마다 $[l_i,u_i]$로 자르기)를 써서 사영된 경사 경로 $x(t)=P(x_k-tg_k)$를 따라 이차 모형 $m_k(x)=f_k+g_k^{\mathsf T}(x-x_k)+(x-x_k)^{\mathsf T}B_k(x-x_k)/2$을 최소화한다. 경로는 성분이 경계에 닿는 점마다 꺾이는 조각별 선분이므로, 꺾이는 점을 차례로 지나며 각 조각에서 이차 함수의 최소를 찾는다. 이 점에서 경계에 닿은 변수들을 \term{활성 집합}{active set}으로 정한다. 둘째는 부분공간 최소화이다. 활성 변수를 경계에 고정하고 나머지 자유 변수에 대해 모형을 최소화한 뒤 그 점을 상자 안으로 되돌린다. 셋째는 $x_k$에서 그 점으로 가는 방향의 직선 탐색으로, 걸음은 상자를 벗어나지 않는 길이로 제한되고 \code{dcsrch}가 찾는다. 이 모든 계산에 $B_k$가 필요한데, L-BFGS-B는 이것을 최근 $m$개의 쌍으로 쓴 \term{압축 표현}{compact representation} $B_k=\theta I-W_kM_kW_k^{\mathsf T}$로 다룬다\cite{byrd1994representations}. 여기서 $W_k$는 $n\times2m$, $M_k$는 $2m\times2m$이고 $\theta$는 $y^{\mathsf T}y/(s^{\mathsf T}y)$이다. 그래서 반복마다 $O(mn)$으로 끝난다. 멈춤은 사영된 기울기 $P(x_k-g_k)-x_k$의 최대 성분이 $10^{-5}$ 이하이거나 함수값의 상대 감소가 $\varepsilon\times10^7$ 이하일 때이다.

\section{Newton-CG}
\label{sec:newton-cg}

Newton 방향은 연립방정식 $\nabla^2f(x_k)\,p=-g_k$의 해이다. 이것을 27장의 공액 기울기법으로 풀면 Hessian 자체는 필요 없고 Hessian과 벡터의 곱만 있으면 된다. 곱은 사용자가 주거나 기울기의 차분 $(\nabla f(x_k+hv)-g_k)/h$로 어림한다. 또 정확히 풀 필요도 없다. 최소점에서 멀 때 정확한 Newton 방향은 어차피 근사이므로, 잔차를 기울기에 비례하는 만큼만 줄이고 멈춘다. 이것이 \term{비정확 Newton 방법}{inexact Newton method}이다.

\begin{theorem}[Dembo--Eisenstat--Steihaug]
\label{thm:inexact-newton}
$f$가 \cref{thm:dennis-more}의 가정을 만족하고 $p_k$가 $\|\nabla^2f(x_k)p_k+g_k\|\le\eta_k\|g_k\|$를 만족한다고 하자. $x_0$이 $x_\ast$에 충분히 가깝고 $\eta_k\le\eta<1$이면 $x_{k+1}=x_k+p_k$는 노름 $\|\nabla^2f(x_\ast)(\cdot)\|$에서 $x_\ast$로 선형 수렴하고, $\eta_k\to0$이면 초선형으로, $\eta_k=O(\|g_k\|^q)$($0<q\le1$)이면 차수 $1+q$로 수렴한다.
\end{theorem}

이 정리\cite{dembo1982}의 초선형 부분은 \cref{thm:dennis-more}에서 바로 나온다. $B_k$ 대신 잔차 $r_k=\nabla^2f(x_k)p_k+g_k$를 보면 $g_k=-\nabla^2f(x_k)p_k+r_k$이고, $\|r_k\|\le\eta_k\|g_k\|$와 $\|g_k\|=O(\|p_k\|)$에서 $\|r_k\|=o(\|p_k\|)$이므로 Dennis--Moré의 증명에서 $g_{k+1}=o(\|p_k\|)$가 그대로 성립하기 때문이다. scipy와 stardust의 Newton-CG는 $\eta_k=\min(1/2,\sqrt{\|g_k\|_1})$을 써서($1$-노름으로 잰다) 처음에는 느슨하게 풀다가 최소점에 다가갈수록 정확히 푼다. $\eta_k=O(\|g_k\|^{1/2})$이므로 차수는 $1.5$이다.

공액 기울기법은 행렬이 양정치일 때의 방법이므로, Hessian이 양정치가 아닌 곳에서는 고쳐 써야 한다. 반복 중에 곡률 $d^{\mathsf T}\nabla^2f\,d$가 음수이거나 $0$에 가까운 방향 $d$를 만나면 공액 기울기 반복을 그만두고 그때까지의 해를 쓴다. 첫 반복에서 곡률이 음수인 방향을 만나면 아직 해가 없으므로 가장 가파른 내리막 방향을 쓴다. 이렇게 하면 방향은 언제나 하강 방향이 되지만 그 대가로 음의 곡률이라는 정보를 버린다. Rosenbrock 함수를 $(-1.2,1)$에서 시작하면 Newton-CG의 경로는 Hessian이 양정치가 아닌 곳을 $42$번 지나고, $83$번 반복한다. $\|g\|_\infty\le10^{-5}$까지 돌린 Newton 방법이 $21$번인 것과 크게 다르다. 반대로 차원이 크면 Newton-CG의 장점이 드러난다. 100차원 Rosenbrock 함수에서 Newton-CG는 $219$번 반복하고 Hessian과 벡터의 곱을 $1{,}775$번 계산해 BFGS의 $538$번보다 적게 반복한다. 다만 걸음이 작아져 멈춘 곳의 함수값이 $2.1\times10^{-5}$로, BFGS의 $1.1\times10^{-13}$보다 최소에서 멀다. 음의 곡률을 버리지 않고 쓰는 방법은 다음 장의 신뢰 영역에서 다룬다.

\begin{table}[t]
\begin{sdtable}{colspec={lcccc}, row{1}={font=\small\bfseries}}
\toprule
방법 & 반복 & 함수값 & 기울기 & 최소점까지의 거리 \\
\midrule
경사 하강법(강한 Wolfe) & $7{,}095$ & $56{,}219$ & --- & $2.3\times10^{-5}$ \\
Newton(Wolfe 걸음) & $21$ & --- & --- & $2.6\times10^{-9}$ \\
BFGS & $32$ & $39$ & $39$ & $5.4\times10^{-8}$ \\
L-BFGS-B ($m=10$) & $36$ & $44$ & $44$ & $2.0\times10^{-6}$ \\
Newton-CG & $83$ & $105$ & $105$ & $3.5\times10^{-5}$ \\
\bottomrule
\end{sdtable}
\caption{2차원 Rosenbrock 함수를 $(-1.2,1)$에서 최소화한 결과. 경사 하강법과 Newton은 $\|g\|_\infty\le10^{-5}$까지 직접 돌렸고, 나머지는 scipy 1.18.1의 기본 설정이다. 최소점까지의 거리는 성분의 최대 오차이다. Newton-CG는 Hessian도 $83$번 계산했다.}
\label{tab:rosenbrock}
\end{table}

\cref{tab:rosenbrock}에서 최소점까지의 거리가 방법마다 다른 것은 멈춤 판정이 다르기 때문이다. BFGS의 판정 $\|g\|_\infty\le10^{-5}$는 Hessian의 가장 작은 고유값이 약 $0.4$인 이 문제에서 위치의 오차로 대략 $10^{-5}/0.4$ 이하를 뜻한다. Newton-CG는 걸음의 크기가 작아지면 멈추므로 판정의 뜻이 또 다르다. 42장에서 말했듯이 \code{converged}는 그 방법의 판정을 만족했다는 뜻이다.

\begin{instardust}[\code{docs/design/optimize.md} §3.5]
\code{sd.optimize.minimize}의 기본 방법은 경계가 없으면 \code{'bfgs'}, 있으면 \code{'l-bfgs-b'}이다. 직선 탐색은 \code{csrc/optimize/wolfe.hpp} 하나를 BFGS와 Newton-CG가 함께 쓰며, scipy의 \code{\_line\_search\_wolfe12}를 numpy의 산술 그대로 옮겼다. BFGS는 이 장의 $O(n^2)$ 갱신을 쓰기 때문에 scipy와 반올림만큼 다른 걸음을 가는데, 설계 문서는 이것을 다변수 최소화의 여섯 방법 가운데 scipy와의 비트 동일에서 일부러 벗어난 유일한 곳으로 적는다. L-BFGS-B는 scipy 1.15부터의 C 판(L-BFGS-B 3.0\cite{morales2011lbfgsb})을 계산 순서까지 옮겼다. scipy가 그 C 파일을 clang의 기본 설정으로 빌드해 \code{a*b+c}를 FMA 하나로 계산하므로 이 파일만 같은 축약을 켜고 식을 같은 꼴로 적었으며, Accelerate의 \code{ddot}이 둘째 인자의 16바이트 정렬에 따라 합의 순서를 바꾸므로 작업 공간도 scipy와 같은 위치로 나눈다. 그래야 경계 있는 문제에서 함수를 부르는 점이 scipy와 비트까지 같다. Newton-CG의 Hessian 곱은 사용자의 \code{hess}나 \code{hessp}, 또는 기울기의 앞쪽 차분이다. 사용자 함수는 Python 콜백이고, 큰 선형대수(BFGS의 $Hg$와 갱신, L-BFGS-B의 방향 계산)에서만 GIL을 놓는다.
\end{instardust}

\section{맺음}

직선 탐색 방법은 방향과 걸음을 나누어 정한다. Wolfe 조건은 걸음이 해야 할 일을 최소한으로 정하고 Zoutendijk의 정리는 방향과 기울기 사이의 각이 직각 쪽으로 다가가지 않는 한, 곧 $\cos\theta_k\ge\delta>0$이면 그것만으로 수렴한다는 것을 보장한다. 수렴 속도는 방향이 정하는데 Newton 방향은 이차로 수렴하지만 Hessian이 비싸고 양정치가 아닐 수 있다. 준Newton 방법은 할선 조건과 가장 작은 변화라는 두 원리로 Hessian을 기울기에서 쌓아 가고, Dennis--Moré 정리는 그 근사가 걸음 방향에서만 맞으면 초선형 수렴에 충분하다고 말한다. 이렇게 나눈 두 몫은 서로 돕기도 해서 Wolfe 조건의 곡률 조건이 BFGS의 양정치성을 지켜 준다. L-BFGS는 행렬을 쌍의 목록으로 바꾸어 큰 문제로 넓히고, Newton-CG는 Hessian을 곱셈으로만 쓴다. 그래도 음의 곡률은 약점으로 남는데 직선 탐색은 정해진 방향을 따라가기만 하므로 방향을 정하는 모형이 나쁠 때 대처할 길이 마땅치 않다. 다음 장의 신뢰 영역 방법은 걸음의 길이를 먼저 정하고 방향을 그 안에서 고른다.

이 장의 내용은 Nocedal과 Wright의 책\cite{nocedal2006} 3장, 6장, 7장과 Dennis와 Schnabel의 책\cite{dennis1996}에 자세히 나온다.

\chapterreferences
